% \setlength{\columnsep}{12pt}
% espaço entre colunas %\setlength{\columnseprule}{0.2pt} % linha vertical entre colunas

%-------------------------------------------------------------------------
% Modelo Trabalho PPGE-UFPB/Estatística
% Normas da ABNT
% Provedor: Laboratório de Economia e Modelagem Aplicada -- LEMA
%-------------------------------------------------------------------------

\documentclass[12pt, openright,oneside, a4paper, english, brazil]{report}

% ----------------------------
% Pacotes básicos 
% ----------------------------
% Fonte e codificação (acentuação)
\usepackage{lmodern}

\renewcommand{\sfdefault}{ppl}
\usepackage[T1]{fontenc}	
\usepackage[utf8]{inputenc}	
\usepackage[brazil]{babel}
\usepackage{comment} % para esconder a tabela Org do compilador
% Referências
%\usepackage[brazilian,hyperpageref]{}	
%\usepackage[alf]{abntex2cite}			
\usepackage[colon]{natbib}
\bibliographystyle{apalike}

%Texto aleatório
\usepackage{lipsum} 

% Tabelas e Figuras
\usepackage{ctable}             
\usepackage{multirow}
\usepackage{array}
\usepackage{float}              
\usepackage{longtable}          
\usepackage{xcolor,colortbl}    
\usepackage{booktabs}           
\usepackage{graphicx}			
\usepackage{subfig}             
\usepackage{caption}            
\usepackage{epstopdf}           

% Equações e símbolos
\usepackage{amsmath}            
\usepackage{nomencl} 
\usepackage{hyperref, cleveref}  

% Texto
\usepackage{indentfirst}		
\usepackage{color}				
\usepackage{microtype} 			
\usepackage{lastpage}			
% ----------------------------------------------------------
% Configurações do texto
% ----------------------------------------------------------
\usepackage{multicol}
\usepackage{array}
\newcolumntype{L}[1]{>{\raggedright\let\newline\\\arraybackslash\hspace{0pt}}m{#1}}
\newcolumntype{C}[1]{>{\centering\let\newline\\\arraybackslash\hspace{0pt}}m{#1}}
\newcolumntype{R}[1]{>{\raggedleft\let\newline\\\arraybackslash\hspace{0pt}}m{#1}}

\makeatletter
%\hypersetup{pdftitle={\@title}, pdfauthor={\@author}}
\makeatother

\def\preambulo{Texto de fixação de conceitos elabotrado por \autor}
%--------------------------------------------------------------

\setlength{\parindent}{1.25cm}
\setlength{\parskip}{0.1cm} 

\usepackage{titlesec}
\usepackage{lipsum}

\makeindex




\title{Conceitos Contabeis}
\def\autor{Oscar Sérgio Monteiro}
%\def\disciplina{SC329 - Sustentabilidade  Ambiental e Social, turma b}
% \def\professor{Elcídio Henriques  Quiraque}
\begin{document}
\makeatletter
\begin{titlepage}
    % \begin{figure}[H]
    % \vspace{-3cm}
    % \centering
    % \includegraphics[height=3cm,keepaspectratio]{config/coverpage}\hfill
    % \par\vspace{1cm}
    % \end{figure}
    % \begin{figure}[H]
    % \vspace{-2cm}
    % \centering
    % \includegraphics[height=3cm,keepaspectratio]{config/contabeis.png}
    % \par\vspace{1cm}
    % \end{figure}
    \vspace{2cm}
	\begin{center}
	    {\huge \bfseries \@title}
	    \begin{flushright}
		\begin{minipage}{8cm}
			\preambulo
		\end{minipage}
	\end{flushright}
	\vspace{3cm}
	    \normalsize
	    \autor
	\end{center}
	
\vfill
\centering
{\large Curitiba-PR, \the\year}

\end{titlepage}
\makeatother


\pdfbookmark[0]{\contentsname}{toc}
\tableofcontents
\newpage


\titleformat{\chapter}[display]{\normalfont\bfseries}{}{0pt}{\Huge  \thechapter. }
\newcommand{\doubletoprule}{%
  \specialrule{.1em}{2pt}{2pt}%
  \specialrule{.1em}{0pt}{0pt}%
}

\chapter{Introdução}
\section{Objeto da contabilidade}
O objeto da contabilidade e o patrimônio das entidades. O contador cuida do patrimônio das entidades.}
\section{Principios baśicos}
são ses os princípios básicos da contabilidade:
\begin{itemize}
\item Entidade\\
  O dono do patrimônio. Visto todo o patrimônio que é objeto do trabalho do contador esta relacionado a uma Entidade. Não se pode relacionar um patrimônio a mais de uma entidade. 
\item Continuidade
  Embora a entidade de um tabalho contábil possa estar relacionado a uma pessoa , o trabalho do contador deve considerar a peretalidade do patrimônio. O trabalho da contabildade trata o tempo como uma repetição dontínua de periodos. 
\item Competência
  por ser uma seguencia de eventos fatos contábeis infinitos, os eventos são agrupados em competências, que é a referênci temporal do tefa.
  Os 
\item Valor Histórico
  A contabilidade se relaciona um processo de analise contábil, a ser visto mais para a frente, mas que se relaciona ao conceito de transação de fato , ou evento que afeta o patrimônio da entidade , assim o pratrimônio é constituido e alterado por transação. Sobre a transação é um evento ou fato que ocorre em um determinad tempo e com um detarminado valor. Se uma fato ou evento não tem um valor definido , nõ tem interesse para a contabilidade.O princípo do Valor histórico estabeleçe, que os eventos devem ser observador pelo valor no instante que ocorreram. NO momento em que o evento afetou o patrim^monio da entidade. Vem sempre o valor é definido de forma explicita.\\
\item Oportunidade
  Esse princípio tem relação com a qualidade da informação contábil. Como o reposte de uma notícia, tem o tempo certo de ser publicada, e não pode ser pela metade ou mal contada.\\
\item Prudência.
  O Padrimônio das entidade faz parte da sua excencia. De forma que a sua a relevêmcia do patrimônio para a entidade deve ser tratada como muita prudência e conservadorismo. O contsdor trbaho em uma entidade e para uma entidade mas deve se manter afastado dela par realizar o seu tablho. Observ a entidade com prdência. Ao medir o patrimônio de uma entidade não pode o contdor criar expectativas de resultado. Deve apenas medir os resultados. Assim deve atuar de forma cauteloza, sem expectativas. Como umm médico que não deve atender aos seus familiares, em função do envolvimento pessoal, que pode fazer com que o médico não vaorize sistomas ue possam levar a um diagnóstico ruim. O contador deve se distanciar da entidade par da nmesma forma não fazer avaliações equivocadas.
\end{itemize}

Em resumo numa sequencia logica.\\
Patrimônio é o objeto da contbilidade. O patrimônio se relaciona a uma \textbf{entidade}. O patrimônio vai de constituindo de fatos contábeis que ocorren no tempo, na ´vida´ da entidade. Assim decorre o principio da \textbf{continuidade} que resalta a perenidade da entidade. E em relação ao instante em que os eventos ocorrem decorre o princípio da \textbf{competência} que é  atributo do fato no qual são agrupados os fatos par a apresentaçao de resultados. O princípio que relaciona valor do fato no tempo é o princípio do \textbf{ valor histórico} que orenta para o registro do fato pelo valor no instante em que ocorreu. Outro princípio relacionado ao ytempo é  princípio da \textbf{oporortunidade}, este principio preconizaque os eventos contábeis , assim como as notícias tem o memento oportuo par a sua divulgação, ademais da necessidade de que o reporte se faça de forma completa e fiel aos fatos. E por tratarmos de patrimônio, e sobre decisões que são tomadas sobre os relalórios contabeis, decorre o princípio da \textbf{prudência} de recomeneta que as analises sejam realizadas de forma conservadora.\\
\end{itemize}
Não obstante não seja tratado como um princípio contábil, uma caracteristica que o registro contábil deve oforecer é o da essenciq sobre a forma. Independente de como o fato contábil se apresente o registro deve ser realizado pela essencia do que o fato representa.\\
\section{caracteristicas qualitativas das informações contábeis}
Podem ser classificadas em caracteristicas fundamentais:
\begin{itemize}
\item relavancia
  pare ser reconhecido como um fato passivel de registro o fato deve ser relevante.\\
\item referencia fidedígna
  relacionada ao princípio da oportunidade o registro contábil deve refletis de forma completa e verdadeira o fato cont´bi\\l
\end{itemize}
Qualidade adicionais da informação contábil:
\begin{itemize}
\item Tempestiva
\item Verificavel
\item Comparável
\item Compreencível
\end{itemize}
Restriçoes e equilíbrio.\\
As informação não podem ser tão detalhadas a ponto de projudicar a sua tempestividade, de outro lado não pode ser tão rápida que possa comprometer sua complitudo. Tamem o custo da sua optenção deve ser equilibrado em relação a sua ree vância.
\section{Princípios da ética profissional do contador - NBC PG 100 (R1)}
Estas normas descrevem os princípios fundamentais de ética para os profissionais da contabilidade, refletindo o reconhecimento da profissão de sua responsabilidade de interesse público. Esses princípios estabelecem o padrão de comportamento esperado do profissional da contabilidade. Os princípios fundamentais são: integridade, objetividade,competência profissional e devido zelo, confidencialidade e comportamento profissionalexit
% \bibliographystyle{apalike}
% \bibliography{referencias}                                                     
 
\end{document}
